% !TEX TS-program = lualatex
% !TEX encoding = UTF-8 Unicode
\documentclass[12pt,a4paper]{article}

\IfFileExists{src/libs/body.tex}{% --- body.tex ---
\usepackage{pdfpages}
\usepackage{fontspec}
% Prefer Times New Roman when available; fall back in container environments.
\IfFontExistsTF{Times New Roman}{
  \setmainfont{Times New Roman}
}{
  \setmainfont{Liberation Serif}
}
\usepackage[vietnamese]{babel}

% Page layout and spacing
\usepackage[left=3cm,right=2cm,top=2cm,bottom=2cm]{geometry}
\usepackage{titlesec}
\usepackage{setspace}

% Silence noisy warnings
\usepackage{silence}
\WarningFilter{transparent}{Loading aborted}
\WarningFilter{hyperref}{The anchor of a bookmark and its parent's}

% Core packages
\usepackage{float}
\usepackage{svg}
\usepackage{graphicx}
\graphicspath{{assets/}{../assets/}{../../assets/}}
\svgpath{{assets/}{../assets/}{../../assets/}}
\usepackage{enumitem}
\usepackage{array}
\usepackage{tabularx}
\usepackage{multirow}
\usepackage[table]{xcolor}
\usepackage{ulem}
\usepackage{xurl}

% ====== THÊM 2 GÓI NÀY ĐỂ FIX LỖI ADJUSTBOX VÀ FOREST ======
\usepackage{adjustbox}
\usepackage[edges]{forest}
% ============================================================

% TikZ
\usepackage{tikz}

% ====== BỔ SUNG FIT, BACKGROUNDS, TREES VÀO DANH SÁCH ======
\usetikzlibrary{calc, arrows.meta, shapes.geometric, positioning, fit, backgrounds, trees}

% Table tuning
\renewcommand{\tabularxcolumn}[1]{m{#1}}
\hbadness=10000
\hfuzz=10000pt

% Shared commands
\newcommand{\subsecheading}[1]{%
    \par\vspace{2pt}%
    \hspace{1.5em}\textbf{#1}%
    \par\vspace{2pt}%
}
\newcommand{\introtext}[1]{%
    \par\hspace{3.0em}#1\par%
}
\newcommand{\StandaloneDocumentSetup}[2]{%
  \pagenumbering{arabic}%
  \ApplyStandardFooter{#1}{#2}%
}
\newcommand{\StandaloneSectionSetup}[2]{%
  \ApplyStandardFooter{#1}{#2}%
}

% Math
\usepackage{amsmath}

% Hyperref
\usepackage{hyperref}
\hypersetup{
    colorlinks=true,
    linkcolor=black,
    filecolor=magenta,
    urlcolor=blue,
}}{% --- body.tex ---
\usepackage{pdfpages}
\usepackage{fontspec}
% Prefer Times New Roman when available; fall back in container environments.
\IfFontExistsTF{Times New Roman}{
  \setmainfont{Times New Roman}
}{
  \setmainfont{Liberation Serif}
}
\usepackage[vietnamese]{babel}

% Page layout and spacing
\usepackage[left=3cm,right=2cm,top=2cm,bottom=2cm]{geometry}
\usepackage{titlesec}
\usepackage{setspace}

% Silence noisy warnings
\usepackage{silence}
\WarningFilter{transparent}{Loading aborted}
\WarningFilter{hyperref}{The anchor of a bookmark and its parent's}

% Core packages
\usepackage{float}
\usepackage{svg}
\usepackage{graphicx}
\graphicspath{{assets/}{../assets/}{../../assets/}}
\svgpath{{assets/}{../assets/}{../../assets/}}
\usepackage{enumitem}
\usepackage{array}
\usepackage{tabularx}
\usepackage{multirow}
\usepackage[table]{xcolor}
\usepackage{ulem}
\usepackage{xurl}

% ====== THÊM 2 GÓI NÀY ĐỂ FIX LỖI ADJUSTBOX VÀ FOREST ======
\usepackage{adjustbox}
\usepackage[edges]{forest}
% ============================================================

% TikZ
\usepackage{tikz}

% ====== BỔ SUNG FIT, BACKGROUNDS, TREES VÀO DANH SÁCH ======
\usetikzlibrary{calc, arrows.meta, shapes.geometric, positioning, fit, backgrounds, trees}

% Table tuning
\renewcommand{\tabularxcolumn}[1]{m{#1}}
\hbadness=10000
\hfuzz=10000pt

% Shared commands
\newcommand{\subsecheading}[1]{%
    \par\vspace{2pt}%
    \hspace{1.5em}\textbf{#1}%
    \par\vspace{2pt}%
}
\newcommand{\introtext}[1]{%
    \par\hspace{3.0em}#1\par%
}
\newcommand{\StandaloneDocumentSetup}[2]{%
  \pagenumbering{arabic}%
  \ApplyStandardFooter{#1}{#2}%
}
\newcommand{\StandaloneSectionSetup}[2]{%
  \ApplyStandardFooter{#1}{#2}%
}

% Math
\usepackage{amsmath}

% Hyperref
\usepackage{hyperref}
\hypersetup{
    colorlinks=true,
    linkcolor=black,
    filecolor=magenta,
    urlcolor=blue,
}}
\IfFileExists{src/libs/header.tex}{% --- header.tex ---
\usepackage{fancyhdr}
\setlength{\headheight}{15.5pt} % Thêm dòng này để fix cảnh báo chiều cao
%\setlength{\headheight}{14.5pt}
\addtolength{\topmargin}{-2.5pt}

\pagestyle{fancy}
\fancyhf{}
\renewcommand{\headrulewidth}{0pt}
\renewcommand{\footrulewidth}{0pt}

\newcommand{\ApplyStandardHeader}{%
  \fancyhead[L]{%
    \begin{minipage}[t]{0.795\linewidth}%
      UIT \textbar{} SS004.F21 \textbar{} Nhóm 03 \textbar{} Đồ án cuối kỳ%
    \end{minipage}%
    \vrule%
    \begin{minipage}[t]{0.185\linewidth}%
      \centering cTetris%
    \end{minipage}%
  }%
}

\ApplyStandardHeader
}{% --- header.tex ---
\usepackage{fancyhdr}
\setlength{\headheight}{15.5pt} % Thêm dòng này để fix cảnh báo chiều cao
%\setlength{\headheight}{14.5pt}
\addtolength{\topmargin}{-2.5pt}

\pagestyle{fancy}
\fancyhf{}
\renewcommand{\headrulewidth}{0pt}
\renewcommand{\footrulewidth}{0pt}

\newcommand{\ApplyStandardHeader}{%
  \fancyhead[L]{%
    \begin{minipage}[t]{0.795\linewidth}%
      UIT \textbar{} SS004.F21 \textbar{} Nhóm 03 \textbar{} Đồ án cuối kỳ%
    \end{minipage}%
    \vrule%
    \begin{minipage}[t]{0.185\linewidth}%
      \centering cTetris%
    \end{minipage}%
  }%
}

\ApplyStandardHeader
}
\IfFileExists{src/libs/footer.tex}{% --- footer.tex ---
\usepackage{lastpage}

\newcommand{\ApplyStandardFooter}[2]{%
  \fancyfoot[L]{%
    \begin{minipage}[t]{0.795\linewidth}%
      Document ID: #1%
    \end{minipage}%
    \vrule%
    \begin{minipage}[t]{0.185\linewidth}%
      \centering\thepage/\pageref{#2}%
    \end{minipage}%
  }%
  \fancypagestyle{plain}{%
    \fancyhf{}%
    \renewcommand{\headrulewidth}{0pt}%
    \renewcommand{\footrulewidth}{0pt}%
    \ApplyStandardHeader%
    \fancyfoot[L]{%
      \begin{minipage}[t]{0.795\linewidth}%
        Document ID: #1%
      \end{minipage}%
      \vrule%
      \begin{minipage}[t]{0.185\linewidth}%
        \centering\thepage/\pageref{#2}%
      \end{minipage}%
    }%
  }%
  \pagestyle{fancy}%
}
}{% --- footer.tex ---
\usepackage{lastpage}

\newcommand{\ApplyStandardFooter}[2]{%
  \fancyfoot[L]{%
    \begin{minipage}[t]{0.795\linewidth}%
      Document ID: #1%
    \end{minipage}%
    \vrule%
    \begin{minipage}[t]{0.185\linewidth}%
      \centering\thepage/\pageref{#2}%
    \end{minipage}%
  }%
  \fancypagestyle{plain}{%
    \fancyhf{}%
    \renewcommand{\headrulewidth}{0pt}%
    \renewcommand{\footrulewidth}{0pt}%
    \ApplyStandardHeader%
    \fancyfoot[L]{%
      \begin{minipage}[t]{0.795\linewidth}%
        Document ID: #1%
      \end{minipage}%
      \vrule%
      \begin{minipage}[t]{0.185\linewidth}%
        \centering\thepage/\pageref{#2}%
      \end{minipage}%
    }%
  }%
  \pagestyle{fancy}%
}
}
\usepackage{docmute}
\hypersetup{pdftitle={Bien ban hop - ctetris}}

\begin{document}
\providecommand{\StandaloneSectionSetup}[2]{\ApplyStandardFooter{#1}{#2}}
\StandaloneSectionSetup{03-meetingMinute.tex}{LastPageMeetingMinute}

\setlist[itemize,1]{label=\textbullet, leftmargin=1.5em, itemsep=0pt, parsep=0pt, topsep=2pt}
\setlist[itemize,2]{label=$\circ$, leftmargin=1.8em, itemsep=0pt, parsep=0pt, topsep=2pt}

\setlength{\parindent}{0pt}
\setlength{\parskip}{0pt}
\singlespacing
% ----------------------------------------------------
% ====================================================
% ===== BIÊN BẢN HỌP NGÀY 12/04/2026 =================
% ====================================================
% ===== TITLE =====
\begin{center}
    {\large\bfseries BIÊN BẢN CUỘC HỌP}
\end{center}
\begin{center}
    \textbf{Ngày: 12 / 04 / 2026}
\end{center}
\vspace{0.4em}

% ===== SECTION 1 =====
\textbf{1. Thông tin chung}
\begin{itemize}
    \item \textbf{Thời gian:} 13h:30 – 13h:55, Chủ nhật, ngày 12/04/2026.
    \item \textbf{Địa điểm:} Trực tuyến qua Google Meet.
    \item \textbf{Tổ chức cuộc họp:} Phạm Tấn Phúc.
    \item \textbf{Thành phần tham dự:} Phạm Tấn Phúc, Ngô Triều Vĩ, Nguyễn Trung Kiên.
    \item \textbf{Mục đích cuộc họp:}
    \begin{itemize}
        \item Review cấu trúc mã nguồn dự án CTetris trên GitHub.
        \item Kiểm tra tài liệu báo cáo trên Latex
        \item Chia sẻ và thảo luận về Gameplay của trò chơi.
    \end{itemize}
\end{itemize}

% ===== SECTION 2 =====
\vspace{4pt}
\textbf{2. Nội dung thảo luận chính}

\subsecheading{2.1. Dự án CTetris (Console Tetris)}
\begin{itemize}
    \item Cấu trúc thư mục mã nguồn: Nhóm đã thống nhất chia dự án thành các module rõ ràng:
    \begin{itemize}
        \item \texttt{app/}: Chứa mã nguồn chính (\texttt{main.cpp}).
        \item \texttt{src/}: Bao gồm các thành phần giao diện console, logic game (gameCore), và hình học (geometry).
        \item \texttt{include/}: Chứa các file header tương ứng cho từng module.
    \end{itemize}
    \item Tài liệu đi kèm: File \texttt{appTree.md} đã liệt kê chi tiết sơ đồ cây của ứng dụng để các thành viên dễ theo dõi.
    \item Tiến độ GitHub: Đã thực hiện cập nhật nội dung game guide.
\end{itemize}

\subsecheading{2.2. Tài liệu Báo cáo (Overleaf/LaTeX)}
\begin{itemize}
    \item Nội dung báo cáo: Nhóm đang hoàn thiện báo cáo.
    \item Cấu trúc file LaTeX: Đã phân chia thành các file nhỏ trong thư mục \texttt{doc/src/content/} (Preface, References, Appendix, GameRequirements...).
    \item Chữ ký xác nhận: Nhóm đã chuẩn bị trang ký tên với chữ ký số của các thành viên.
\end{itemize}

\subsecheading{2.3. Dự án cTetris}
\begin{itemize}
    \item Gameplay: Trò chơi cTetris
    \item Cơ chế điều khiển:
    \begin{itemize}
        \item Phím A/D hoặc Mũi tên để di chuyển.
        \item Phím W hoặc Mũi tên lên xoay.
    \end{itemize}
    \item Quy tắc tính điểm: đạt 1 dòng được +1 điểm. Tốc độ game tăng dần theo điểm số.
\end{itemize}

\subsecheading{2.4. Thảo luận khác}
\begin{itemize}
    \item Nhóm đã kiểm tra các kênh trao đổi trên Slack (\texttt{\#all-SS004f21}) để cập nhật thông báo từ giảng viên.
    \item Thống nhất hạn chót hoàn thiện các phần còn lại là trước 20:00 cùng ngày.
\end{itemize}

% ===== SECTION 3 =====
\vspace{2pt}
\textbf{3. Kết luận và Phân công tiếp theo:}
\begin{itemize}
    \item \textbf{Phạm Tấn Phúc:} Tổng hợp mã nguồn cuối cùng và kiểm tra tính ổn định của game.
    \item \textbf{Tạ Minh Trường:} Hoàn thiện nốt các phần nội dung còn trống trong file báo cáo LaTeX.
    \item \textbf{Cả nhóm:} Cùng rà soát lại file PDF xuất ra từ Overleaf để đảm bảo không lỗi định dạng trước khi nộp.
\end{itemize}

% ====================================================
% ===== BIÊN BẢN HỌP NGÀY 15/04/2026 =================
% ====================================================
\newpage

\begin{center}
    {\large\bfseries BIÊN BẢN CUỘC HỌP}
\end{center}
\begin{center}
    \textbf{Ngày: 15 / 04 / 2026}
\end{center}
\vspace{0.4em}

% ===== SECTION 1 =====
\textbf{1. Thông tin chung}
\begin{itemize}
    \item \textbf{Thời gian:} 20:20 - 20:35, Thứ Tư, ngày 15/04/2026.
    \item \textbf{Địa điểm:} Trực tuyến qua Google Meet.
    \item \textbf{Tổ chức cuộc họp:} Phạm Tấn Phúc.
    \item \textbf{Thành phần tham dự:} Phạm Tấn Phúc, Ngô Triều Vĩ, Tạ Minh Trường, Đàm Gia Thái, Nguyễn Trung Kiên.
    \item \textbf{Mục đích cuộc họp:}
    \begin{itemize}
        \item Kiểm tra tiến độ công việc trên bảng quản trị Trello.
        \item Review cấu trúc mã nguồn dự án trên GitHub và môi trường lập trình Codespaces.
        \item Thống nhất các đầu việc cần hoàn thiện cho báo cáo cuối khóa.
    \end{itemize}
\end{itemize}

% ===== SECTION 2 =====
\vspace{4pt}
\textbf{2. Nội dung thảo luận chính}

\subsecheading{2.1. Quản lý công việc (Trello)}
\begin{itemize}
    \item Nhóm đã tiến hành rà soát các thẻ công việc (cards) trên Trello của từng thành viên:
    \begin{itemize}
        \item \textbf{Thái:} Đã hoàn thành các phần việc được giao, chuyển thẻ sang cột "Done".
        \item \textbf{Trường:} Đang thực hiện viết User Acceptance Test (UAT) cho module gameCore.
        \item \textbf{Vĩ:} Đã cập nhật các đầu việc liên quan đến giao diện console.
        \item \textbf{Phúc:} Thực hiện điều phối và di chuyển các task đã hoàn thành để theo dõi tiến độ tổng thể.
    \end{itemize}
\end{itemize}

\subsecheading{2.2. Quản lý mã nguồn và Tài liệu (GitHub \& Overleaf)}
\begin{itemize}
    \item \textbf{GitHub:} Review cấu trúc thư mục dự án ctetris:
    \begin{itemize}
        \item \texttt{app/}: Chứa logic chính và file \texttt{main.cpp}.
        \item \texttt{doc/}: Chứa file \texttt{appTree.md} (sơ đồ cây ứng dụng) và tài liệu \texttt{docTree.md}.
        \item \texttt{src/}: Chia nhỏ thành các module gameConsole, gameCore, gameStory, geometry.
    \end{itemize}
    \item \textbf{Overleaf (LaTeX):} Nhóm đang soạn thảo báo cáo "Kỹ năng nghề nghiệp - SS004.F21.CNTT".
    \begin{itemize}
        \item Đã hoàn thiện trang bìa, mục lục và các phần nội dung kỹ thuật (SRS, SDS).
        \item Đã có trang chữ ký xác nhận của các thành viên trong nhóm.
    \end{itemize}
\end{itemize}

\subsecheading{2.3. Thảo luận về nội dung báo cáo}
\begin{itemize}
    \item Nhóm đã xem xét các file nội dung trong \texttt{doc/src/content/}:
    \begin{itemize}
        \item \texttt{01-Preface.tex}: Lời mở đầu.
        \item \texttt{02-gameStoryRequirements.tex}: Yêu cầu cho phần cốt truyện game.
        \item \texttt{09-gameConsoleGuide.tex}: Hướng dẫn sử dụng giao diện console.
    \end{itemize}
    \item Thảo luận về cách trình bày sơ đồ máy trạng thái (Finite State Machine) và kiến trúc hệ thống trong báo cáo.
\end{itemize}

\subsecheading{2.4. Công tác chuẩn bị khác}
\begin{itemize}
    \item Nhóm đã kiểm tra lại các thông báo trên Slack và trao đổi qua trello để đảm bảo mọi thành viên nắm bắt được lịch họp và hạn nộp bài.
\end{itemize}

% ===== SECTION 3 =====
\vspace{2pt}
\textbf{3. Kết luận và Giao việc:}
\begin{itemize}
    \item \textbf{Trường:} Hoàn thành tài liệu UAT và cập nhật vào file báo cáo.
    \item \textbf{Cả nhóm:} Kiểm tra lại định dạng file PDF xuất ra từ Overleaf, đặc biệt là phần hình ảnh và sơ đồ.
    \item \textbf{Hạn chót:} Hoàn thiện toàn bộ báo cáo và mã nguồn để chuẩn bị nộp bài đúng hạn.
\end{itemize}

% ====================================================
% ===== BIÊN BẢN HỌP NGÀY 06/05/2026 =================
% ====================================================
\newpage

\begin{center}
    {\large\bfseries BIÊN BẢN CUỘC HỌP}
\end{center}
\begin{center}
    \textbf{Ngày: 06 / 05 / 2026}
\end{center}
\vspace{0.4em}

% ===== SECTION 1 =====
\textbf{1. Thông tin chung}
\begin{itemize}
    \item \textbf{Thời gian:} Thứ Sáu, ngày 06/05/2026.
    \item \textbf{Địa điểm:} Họp trực tuyến qua Google Meet.
    \item \textbf{Tổ chức cuộc họp:} Phạm Tấn Phúc.
    \item \textbf{Thành phần tham dự:} Phạm Tấn Phúc, Nguyễn Trung Kiên, Đàm Gia Thái, Ngô Triều Vĩ, Tạ Minh Trường.
    \item \textbf{Mục đích cuộc họp:}
    \begin{itemize}
        \item Họp tổng kết nhóm.
        \item Cập nhật tiến độ công việc trên Trello cho các Tuần 1, 2 và 3.
        \item Kiểm tra tình trạng Pull Request (PR) trên GitHub và phân công công việc giai đoạn cuối.
    \end{itemize}
\end{itemize}

% ===== SECTION 2 =====
\vspace{4pt}
\textbf{2. Nội dung thảo luận chính}

\subsecheading{2.1. Báo cáo tiến độ trên Trello}
\begin{itemize}
    \item \textbf{Tuần 1:} Đã hoàn thành toàn bộ các công việc (Version 1).
    \item \textbf{Tuần 2:} Các thành viên tiếp tục thực hiện task đã được phân công trên Trello. \textbf{Lưu ý quan trọng từ Giảng viên:} Yêu cầu toàn bộ code phải được viết trên một file duy nhất. Phạm Tấn Phúc đã tiến hành revert lại code để khớp với yêu cầu của tài liệu.
    \item \textbf{Tuần 3:} Đã thiết lập xong các task trên Trello. Trọng tâm của tuần 3 là thực hành cố tình tạo ra các lỗi (size, conflict) để các thành viên trải nghiệm cách xử lý và giải quyết xung đột code.
\end{itemize}

\subsecheading{2.2. Quản lý mã nguồn trên GitHub}
\begin{itemize}
    \item Việc phân công theo từng nhánh (branch) cho các thành viên đã hoàn tất.
    \item Tình trạng Pull Request (PR) tuần 2 (được Phạm Tấn Phúc kiểm tra trực tiếp trong cuộc họp):
    \begin{itemize}
        \item PR của Ngô Triều Vĩ: Đã Pass (Thành công).
        \item PR của Nguyễn Trung Kiên: Đang gặp trạng thái Conflict (Xung đột).
        \item PR của Phạm Tấn Phúc: Tạm dừng phát triển v3 cũng như sửa lỗi v2, chuyển codebase về lại 1 file main.cpp như gợi ý của thầy.
    \end{itemize}
\end{itemize}

\subsecheading{2.3. Tài liệu đồ án}
\begin{itemize}
    \item Bước vào tuần cuối, các thành viên cần tập trung thời gian để hoàn thiện các phần tài liệu đồ án còn lại.
\end{itemize}

\subsecheading{2.4. Ý kiến các thành viên}
\begin{itemize}
    \item \textbf{Tạ Minh Trường:} Nhất trí thực hiện công việc theo đúng phân công trên Trello, không có câu hỏi thêm.
    \item \textbf{Ngô Triều Vĩ:} Không có thắc mắc hay bổ sung gì thêm.
\end{itemize}

% ===== SECTION 3 =====
\vspace{2pt}
\textbf{3. Kết luận và Phân công tiếp theo:}
\begin{itemize}
    \item \textbf{Phạm Tấn Phúc:} Kiểm tra, giải quyết lỗi Conflict trong PR của Nguyễn Trung Kiên và thực hiện merge code.
    \item \textbf{Toàn bộ thành viên:} Bám sát tiến độ và hoàn thành các task Tuần 2 \& Tuần 3 trên Trello. Viết và hoàn thiện tài liệu đồ án cho tuần cuối.
\end{itemize}

\label{LastPageMeetingMinute} 
\end{document}